\section{Fine-Tuning 的类型与适用场景}

\subsection{三种 Fine-Tuning 类型}

讲者将 Fine-Tuning 分为三类，各自对应不同的数据规模需求和目标：

\begin{table}[H]
\centering
\caption{Fine-Tuning 类型对比}
\begin{tabular}{lccc}
\toprule
\textbf{类型} & \textbf{目标} & \textbf{数据规模} & \textbf{示例} \\
\midrule
通用型 (General Purpose) & 创建全能 chatbot & $>$100万样本 & ChatGPT, Claude \\
领域型 (Domain Specific) & 嵌入领域知识 & 数千$\sim$数十万 & 医疗/法律/金融 \\
任务型 (Task Specific) & 学习单一功能 & 数百$\sim$数千 & 摘要、拼写检查 \\
\bottomrule
\end{tabular}
\end{table}

\begin{warningbox}{数据需求受多因素影响}
上述样本数量仅为粗略估计。实际需求取决于：(1) \textbf{模型规模}——模型越大，所需样本越少；(2) \textbf{任务复杂度}——若 base model 对目标语言/领域的知识接近于零（如从未见过的小语种），fine-tuning 的效果会非常有限，因为缺乏足够 token 来"从头学习"。
\end{warningbox}

\subsection{何时应该 Fine-Tune}

讲者建议遵循一个明确的决策流程：

\begin{enumerate}
    \item \textbf{优先尝试 In-Context Learning / RAG}：门槛最低，便于快速验证。
    \item 若效果不满意，考虑以下四种 fine-tuning 动机：
    \begin{itemize}
        \item \textbf{改变输出风格/格式}：如调整模型生成邮件的语气
        \item \textbf{添加领域知识}（浅层）：如教模型关于你公司的事实
        \item \textbf{知识蒸馏}：将 GPT-4 级别能力蒸馏到更小的模型中，降低成本和延迟
        \item \textbf{提升特定任务质量}：如 diagram 生成，可在窄任务上超越 frontier model
    \end{itemize}
\end{enumerate}

\begin{knowledgebox}{企业选择 Fine-Tuning 的商业动机}
除了技术层面，企业选择开源模型 + fine-tuning 路线还有两个重要的非技术原因：(1) \textbf{数据控制}——不希望将敏感数据通过 API 发送到云端；(2) \textbf{定制化}——希望深度定制模型的行为、语气和输出格式，创建专属模型。
\end{knowledgebox}

\subsection{本章小结}

Fine-Tuning 分为通用型、领域型和任务型三个层次，数据需求依次递减。实际操作中应先尝试 in-context learning，确认不足后再考虑 fine-tuning。Fine-Tuning 的价值不仅在于技术提升，还在于给予企业对模型和数据的控制权。

\newpage
